\documentclass[10pt,a4paper]{amsart}
\usepackage[margin=19mm]{geometry}
\usepackage{amsmath,amssymb,mathtools}
\usepackage[scr=rsfs,cal=euler]{mathalfa}
\usepackage{xfrac}
\usepackage[unicode,hidelinks,bookmarks=false]{hyperref}
\newcommand{\ie}{i.e.\ }
\newcommand{\cf}{cf.\ }
\newcommand{\resp}{respectively\ }
\newcommand{\Subsection}{\subsection}
\providecommand{\nobreakdash}{\nobreak\hbox{-}}
\setlength{\emergencystretch}{2em}
\multlinegap0pt
\usepackage{amssymb}
\usepackage{bm}
\newtheorem{lemma}{Lemma}
\newtheorem{theorem}{Theorem}
\newcommand{\Fr}{\operatorname{Fr}}
\newcommand{\Mod}{\operatorname{Mod}}
\newcommand{\dom}{\operatorname{dom}}
\newcommand{\inv}{^{-1}}
\newcommand{\spt}{\operatorname{spt}}
\title{Embeddability and rectifiability of Lipschitz differentiability spaces}
\author{Ivan Caamano, Sylvester Eriksson-Bique, Elefterios Soultanis}
\date{Source: arXiv:2609.06700v1 (CC BY 4.0)}
\begin{document}
\maketitle

\noindent\emph{Source and licence.} Embeddability and rectifiability of Lipschitz differentiability spaces. Version arXiv:2609.06700v1 (CC BY 4.0), distributed by arXiv under the Creative Commons Attribution 4.0 International licence, http://creativecommons.org/licenses/by/4.0/, as recorded in the arXiv licence metadata for this exact version and shown on the abstract page https://arxiv.org/abs/2609.06700v1. Full licence notice: \texttt{licenses/arxiv-2609.06700-CC-BY-4.0.txt}.\par\smallskip

\noindent\textit{Editorial scope.} Counted unit: the article's theorem thm:decomp-bundle (source0.tex lines 712--714). The article's complete original proof of it is delivered with it (source0.tex lines 907--943, one \texttt{proof} environment of 37 lines, which the article places away from the statement). No mathematics is written, repaired or re-derived: the statement and the proof are the article's own lines, in the article's own notation, and no retained line is substituted or re-flowed.  Retained uncounted as the source-local context the proof needs: lines 694--699; lines 885--887.  Nothing was omitted from any retained span, so the package carries no omission record.  No retained line contains a citation, so the wrapper carries no bibliography and no bibliography entry.  No retained line contains a figure environment, an image-inclusion command or drawing code, so no image is delivered and the package declares no asset dependency.  Editorial formatting only, in addition: an AMS \texttt{amsart} preamble that restates the article's own declarations and macros used by the retained lines, namely the environments \texttt{lemma}, \texttt{theorem} on separate counters, together with the standard packages those lines need.  The retained environments print in this document as Lemma 1, Theorem 1 in order of appearance, whereas the article prints the same items with the numbering of its own full document.  The complete native source of the article as distributed in the arXiv e-print for this exact version is bundled unchanged as \texttt{sources/209560/source0.tex}.
\begin{lemma}\label{lem:basic-prop-of-decomp-bundle}
The decomposability bundle of $\mu$ (if it exists) is $\mu$-a.e. unique: If $x\mapsto T_\mu(x)$ and $x\mapsto \widetilde T_\mu(x)$  are decomposability bundles as above, then $T_\mu=\widetilde T_\mu$ $\mu$-a.e. on $Y$. Moreover, if $\tilde\mu\ll\mu$ then 
\begin{align*}
    T_{\tilde\mu}(x)=T_\mu(x)\quad\tilde\mu\mbox{-a.e. }\; x.
\end{align*}
\end{lemma}

\begin{lemma}\label{lem:phi-distinguishes-subspaces}
If $V \subset W \subset Y_0$ are two subspaces, and $\phi(W) \le \phi(V)$, then $V = W$.
\end{lemma}

\begin{theorem}\label{thm:decomp-bundle}
Let Y be a Banach space so that $Y$ and $Y^*$ have the Radon-Nikodym property. Then any locally finite Borel regular measure, concentrated on a separable subset, admits a decomposability bundle.
\end{theorem}

\begin{proof}[Proof of Theorem \ref{thm:decomp-bundle}]
Let $\mu$ be a boundedly finite Borel regular measure on $Y$  that is concentrated on a separable set. We can assume that $\mu$ is finite, cf. Lemma \ref{lem:basic-prop-of-decomp-bundle}. Then $\spt\mu$, and thus $Y_0 = \overline{\mbox{span}}\{\spt\mu\}$, is separable.

Consider the collection $\mathcal C$ of Borel bundles $x \mapsto V_x \in \operatorname{Gr}(Y_0)$ with the following property:
\begin{align}\label{eq:tang-bundle}
\gamma_t'\in V_{\gamma(t)}\quad \mbox{a.e.}\; t\in \gamma\inv(\spt\mu)    
\end{align}
for $\Mod_\infty^\mu$-a.e. $\gamma\in \Fr(\spt\mu)$. The collection $\mathcal C$ is non-empty because the constant bundle $x\mapsto Y_0$ belongs to $\mathcal C$. We moreover remark that if $(V^i)\subset \mathcal C$ is a countable set, then the pointwise (countable) intersection $\displaystyle \bigcap_iV^i$ belongs to $\mathcal C$. Indeed, if $\Gamma_i$ are the $\Mod_\infty^\mu$-null families of fragments where \eqref{eq:tang-bundle} fails for $V^i$, then $\Gamma:=\bigcup_i\Gamma_i$ is $\Mod_\infty^\mu$-null and for each $\gamma\notin\Gamma$ there are null-sets $N_i\subset \dom(\gamma)$ so that we have $\gamma_t'\in V_{\gamma(t)}^i$ for $t\in \gamma\inv(\spt\mu)\setminus N_i$.  Thus $\gamma_t'\in \bigcap_iV_{\gamma(t)}^i$ for all $t\in \dom(\gamma)\setminus N$, where $N:=\bigcup_iN_i\subset \dom(\gamma)$. Since $N$ is a null-set, $\bigcap_i V^i$ satisfies \eqref{eq:tang-bundle} and is thus in $\mathcal{C}$.


 Note that, for each $V \in \mathcal C$, and $e\in Y_0^*$ the map $\phi_e:x\mapsto \|e|_{V(x)}\|$ is $\mu-$measurable. Indeed, $\phi_e^{-1}(a,\infty)=\pi(\{(x,v): e(v)\in (a,b), |v|\leq 1, v\in V(x)\})$ is a Suslin set, since $\{(x,v): e(v)\in (a,b), |v|\leq 1, v\in V(x)\}$ is Borel. Here, $\pi:Y\times Y \to Y$ is the projection map $\pi(x,v)=x$. As a sum of such maps $x \mapsto \phi(V_x)$ is $\mu$-measurable.m Consider the minimization problem
\[
\inf_{V \in \mathcal C} \int \phi(V_x)\, d\mu(x) =: A,
\]
and let $(V^n) \subset \mathcal C$ be a minimizing sequence. The intersection $W_x = \bigcap V^n_x$ belongs to $\mathcal C$ as discussed above. It follows
that
\[
A = \int \phi(W_x)\, d\mu = \lim_{n\to\infty} \int \phi(V^n_x)\, d\mu(x) = A,
\]
i.e. $W \in C$ minimizes
\[
\int \phi(V_x)\, d\mu
\]
among $V \in \mathcal C$. We claim that $W$ is $\mu$-a.e. minimal with respect to inclusion, i.e. if $V \in \mathcal C$ then $W_x \subset V_x$ $\mu$–a.e.

Indeed, given $V \in \mathcal C$, the intersection $V \cap W$ belongs to $\mathcal C$. The pointwise inequality
\[
\phi(V_x \cap W_x) \le  \phi(W_x) \quad \mu\text{–a.e.}
\]
and the integral inequality
\[
\int \phi(W_x)\, d\mu \le \int \phi(V_x \cap W_x)\, d\mu
\]
together imply $\phi(W_x) = \phi(V_x \cap W_x)$ $\mu$–a.e., which by Lemma \ref{lem:phi-distinguishes-subspaces} implies that $W = V \cap W$ $\mu$–a.e., completing the proof of the claim. 

We have proven that $W$ is the $\mu$-a.e. minimal Borel bundle satisfying \eqref{eq:tang-bundle}, thus it is the decomposability bundle of $\mu$.
\end{proof}

\end{document}
